\documentclass[11pt]{article}
\usepackage[margin=1in]{geometry}
\usepackage{amsmath,amssymb}
\usepackage{enumitem}
\title{COMP3242 Week 1 Practice Exam (Questions)}
\date{}
\begin{document}
\maketitle

\section*{Instructions}
Answer all questions. Show key steps for computational questions.

\section*{Q1: Tensor Shapes (Concept + Calculation)}
A mini-batch of RGB images has shape $(B, C, H, W) = (16, 3, 32, 32)$.
\begin{enumerate}[label=(\alph*)]
  \item What is the shape after flattening each image to a vector?
  \item A linear layer maps each flattened vector to 100 outputs. What is the weight matrix shape and bias shape?
  \item What is the output batch shape?
\end{enumerate}

\section*{Q2: Matrix-Vector Multiplication}
Let
\[
A=\begin{bmatrix}1&2&0\\-1&3&4\end{bmatrix},\quad x=\begin{bmatrix}2\\-1\\1\end{bmatrix},\quad b=\begin{bmatrix}1\\-2\end{bmatrix}.
\]
Compute $y=Ax+b$.

\section*{Q3: Inner Product and Norm}
Let $u=[1,-2,2]^T$ and $v=[3,0,-1]^T$.
\begin{enumerate}[label=(\alph*)]
  \item Compute $u^Tv$.
  \item Compute $\lVert u\rVert_2$.
\end{enumerate}

\section*{Q4: Broadcasting}
Given tensors $X\in\mathbb{R}^{8\times 4}$ and $b\in\mathbb{R}^{4}$, explain why $X+b$ is valid in PyTorch and state the output shape.

\section*{Q5: PyTorch Debugging}
Consider the code:
\begin{verbatim}
logits = model(x)              # shape (32, 10)
labels = labels.float()        # labels originally class indices
loss = torch.nn.CrossEntropyLoss()(logits, labels)
\end{verbatim}
Identify the bug and provide a correct fix.

\end{document}
